\section{Exponent blocks instead of prime-by-prime logarithms}
\label{sec:ep-start}
\noindent\textbf{Continuation status.} Author: ChatGPT. The preceding moving-support
supplement proves a restricted theorem, not standard $abc$. We retain its
objects and replace only the representation used in its logarithmic forms.
The new results below are ordinary mathematical proofs using two explicitly
stated established inputs. The companion Lean file covers elementary
regrouping and obstruction certificates, not these external estimates or a
closed term of type \texttt{ABCConjecture}. No claim of priority or independent
peer review is made.

Write $T=abc$, $R=\rad(T)$, $t=\log c$, and
\[
 z=a+b\zeta,\qquad \zeta^2-\zeta+1=0,\qquad
 M=N(z)=a^2+ab+b^2.
\]
As before, $a,b$ are positive coprime integers and $c=a+b$.
For $N\ge1$ define
\[
 E_3(N)=\prod_{p\mid N}p^{(v_p(N)-3)_+},\qquad
 C_3(N)=\prod_{p\mid N}p^{(3-v_p(N))_+}.
\]
For a real $Y\ge1$, a superscript $>Y$ restricts either product to primes
$p>Y$. Define also the positive rational, compensated tail
\[
 W_Y(T)=E_3^{>Y}(T)/C_3^{>Y}(T).
\]
It can be less than one. These are functions of actual factorizations, not
abstract capacities.

We use the prime-norm factorization already established in the moving-support
supplement:
\begin{equation}\label{eq:ep-factor}
 z=u(1+\zeta)^{e_0}\prod_{i=1}^r\pi_i^{e_i},\qquad
 M=3^{e_0}\prod_{i=1}^r q_i^{e_i},
\end{equation}
where $e_0\in\{0,1\}$, $q_i=N(\pi_i)\equiv1\pmod3$ are distinct primes,
$e_i\ge1$, and a conjugate pair of factors is never used together.
The only positive triple with $r=0$ is $(1,1,2)$; it is handled directly.
The norm--boundary identity $\gcd(M,T)=1$ and
\begin{equation}\label{eq:ep-cubic}
 z^3-\bar z^3=3\sqrt{-3}\,T,\qquad
 \frac34 c^2\le M\le c^2
\end{equation}
will be used with their full arithmetic hypotheses.

\begin{definition}[A partition into exponent blocks]\label{def:ep-blocks}
Let $\mathcal P=\{I_1,\ldots,I_m\}$ be any partition of $\{1,\ldots,r\}$
into nonempty subsets. Put
\[
 k_j=\gcd\{e_i:i\in I_j\},\quad
 w_j=\prod_{i\in I_j}\pi_i^{e_i/k_j},\quad
 Q_j=N(w_j)=\prod_{i\in I_j}q_i^{e_i/k_j},\quad S=\sum_{j=1}^m k_j.
\]
Thus $k_j\ge1$, $Q_j\ge7$, and
\begin{equation}\label{eq:ep-regroup}
 z=u(1+\zeta)^{e_0}\prod_{j=1}^m w_j^{k_j},\qquad
 \sum_j k_j\log Q_j=\log M-e_0\log3\le2t.
\end{equation}
The one-block partition has $k_1=g:=\gcd(e_1,\ldots,e_r)$.
Grouping equal values of $e_i$ is another allowed partition; its number of
blocks is the number of distinct positive exponents, not the number of primes.
\end{definition}

All assertions in \eqref{eq:ep-regroup} follow from integer divisibility and
associative commutative multiplication. In particular grouping does not create
an algebraic root: the displayed $w_j$ are elements of the original quadratic
ring. Allowing a partition with a few blocks also handles $g=1$, for example
when two blocks have coprime contents $k$ and $k+1$.

\subsection{The two established inputs and their exact use}
We use inequalities (1.2) and the first inequality of Theorem 1.3 of
Bugeaud \cite{EPBugeaud2022}. In a number field of fixed degree $D$, the first
asserts, for a nonzero integer linear form in $s$ fixed determinations of
logarithms,
\[
 \log|b_1\log\alpha_1+\cdots+b_s\log\alpha_s|
 \ge -c_D^s\prod_i h^*(\alpha_i)\log\max(3,|b_1|,\ldots,|b_s|).
\]
Zero coefficients can be deleted and a nonzero last coefficient chosen. The
second, for $s\ge2$ and a product different from one, asserts
\[
 v_{\mathfrak p}\left(\prod_i\alpha_i^{b_i}-1\right)
 <d_D^s p^D\prod_i h^*(\alpha_i)\log\max(3,|b_1|,\ldots,|b_s|).
\]
Here $v_{\mathfrak p}(p)=1$ and $h^*(\alpha)=\max(1,h(\alpha))$.
The constants depend only on $D$. We use neither a conjectural refinement nor
one branch of an unproved alternative. Our applications have $D=2$ throughout.
These results are external established inputs, not new theorems proved or
axiomatized in the companion Lean development.

\begin{theorem}[Blockwise archimedean and non-archimedean bounds]
\label{thm:ep-main}
There is an effective absolute constant $C>2$ such that, for every partition
in Definition~\ref{def:ep-blocks}, the quantity
\begin{equation}\label{eq:ep-penalty}
 \mathcal B_{\mathcal P}(z)=C^{m+1}\log(3+S)\prod_{j=1}^m\log Q_j
\end{equation}
satisfies, with $\lambda=(z/\bar z)^3$,
\begin{align}
 |\lambda-1|&\ge\exp(-\mathcal B_{\mathcal P}(z)),\label{eq:ep-angle}\\
 v_p(T)&\le p^2\mathcal B_{\mathcal P}(z)\quad(p\mid T),\label{eq:ep-adic}\\
 c^3&\le8\exp(\mathcal B_{\mathcal P}(z))R^3
 \frac{E_3(T)}{C_3(T)},\label{eq:ep-height}\\
 \log\prod_{\substack{p\mid T\\p\le Y}}p^{v_p(T)}
 &\le\mathcal B_{\mathcal P}(z)Y^3\log Y\quad(Y\ge2).
 \label{eq:ep-small}
\end{align}
The same conclusions hold with the minimum of
$\mathcal B_{\mathcal P}(z)$ over all partitions. No sublinear bound for this
minimum on all triples is asserted.
\end{theorem}
\begin{proof}
Put $\eta_j=(w_j/\bar w_j)^3$. Both conjugates of $w_j$ have absolute value
$\sqrt{Q_j}$, so $h(w_j)=\tfrac12\log Q_j$ and the height laws give
$h^*(\eta_j)\le3\log Q_j$. None is one: otherwise $w_j^3=\bar w_j^3$,
so each split prime factor of $w_j$ would also divide its conjugate, contrary
to \eqref{eq:ep-factor}. The unit and ramified factors contribute respectively
one and $(-1)^{e_0}$. Consequently
\[
 \lambda=(-1)^{e_0}\prod_j\eta_j^{k_j}\ne1,
\]
where nonvanishing also follows directly from \eqref{eq:ep-cubic} and $T>0$.
Choose principal arguments $\log\eta_j=i\theta_j$, $|\theta_j|\le\pi$, and
an integer $v$ for which
\[
 \Lambda=\sum_j k_j\log\eta_j+(e_0-2v)\log(-1)
\]
has imaginary part in $[-\pi,\pi]$. It is nonzero, and
$|e_0-2v|\le1+S$. Apply the complex bound with at most $m+1$ logarithms,
then use $|e^{ix}-1|\ge2|x|/\pi$ on that interval. All factors $3^m$,
the constants for degree two, and the additive $\log(\pi/2)$ are absorbed
in one absolute $C$, proving \eqref{eq:ep-angle}.

For $p\mid T$, the identity $\gcd(M,T)=1$ implies that $z,\bar z$ and every
$w_j,\bar w_j$ are units at places over $p$. Equation \eqref{eq:ep-cubic}
gives
\[
 v_{\mathfrak p}(\lambda-1)
 =v_p(T)+v_{\mathfrak p}(3\sqrt{-3})\ge v_p(T).
\]
Apply the second external bound to the $m+1\ge2$ numbers
$\eta_1,\ldots,\eta_m,-1$ with coefficients $k_1,\ldots,k_m,e_0$.
Retaining a zero last coefficient is permitted here. This proves
\eqref{eq:ep-adic}, including the primes two and three. Summing and using
$\sum_{p\le Y}p^2\log p\le Y^3\log Y$ proves \eqref{eq:ep-small}.

Finally \eqref{eq:ep-cubic} and $M\ge3c^2/4$ give
$T\ge c^3\exp(-\mathcal B_{\mathcal P})/8$.
Prime by prime $T C_3(T)=R^3E_3(T)$, proving
\eqref{eq:ep-height}. Dropping $C_3(T)\ge1$ recovers the one-sided
bound. The compensated version retains all low-depth credit.
There are finitely many partitions and every one gives all the stated
inequalities. Choosing a minimizing partition is therefore legitimate.
\end{proof}

\section{Removing both prime-size and prime-count restrictions on explicit profiles}
\label{sec:ep-classes}
\begin{corollary}[One block, without fixing the root]\label{cor:ep-content}
Let $g=\gcd(e_1,\ldots,e_r)$. Then an admissible penalty is
\[
 C^2\frac{\log M-e_0\log3}{g}\log(3+g)
 \le2C^2t\frac{\log(3+g)}g.
\]
Neither the size nor the number of norm primes is restricted. For $g\ge64$,
put $Y=g^{1/6}$. The full small-prime part satisfies
\begin{equation}\label{eq:ep-content-small}
 \log\prod_{\substack{p\mid T\\p\le g^{1/6}}}p^{v_p(T)}
 \le\frac{C^2}{3}t\frac{\log(3+g)\log g}{\sqrt g}.
\end{equation}
Both ratios to $t$ tend to zero as $g\longrightarrow\infty$, uniformly in
all other arithmetic data.
\end{corollary}
\begin{proof}
Take the one-block partition in Theorem~\ref{thm:ep-main};
$\log Q_1=(\log M-e_0\log3)/g$. Substitute $Y=g^{1/6}$ into
\eqref{eq:ep-small}. The displayed limits are elementary.
\end{proof}

\begin{corollary}[A uniform restriction on unbalanced norm powers]
\label{cor:ep-unbalanced}
For each $0<\eta<1$ there is an effective $G_\eta$ such that every primitive
positive triple with $g\ge G_\eta$ has
\[
 \min(a,b)>c^{1-\eta}.
\]
In particular this is uniform even when $M=3^{e_0}Q^g$ has arbitrarily large
or numerous split prime factors in $Q$.
\end{corollary}
\begin{proof}
If $\min(a,b)\le c^{1-\eta}$ then $T\le c^{3-\eta}$.
The lower bound for $T$ in the proof of Theorem~\ref{thm:ep-main} yields
\[
 \eta t\le\log8+2C^2t\frac{\log(3+g)}g.
\]
But $t\ge g\log7/2$. Choose $G_\eta$ so large that, for every
$g\ge G_\eta$, each of $2\log8/(g\log7)$ and
$2C^2\log(3+g)/g$ is at most $\eta/3$. Dividing by $t$ gives a
contradiction. Such a threshold is effective and depends only on $\eta$.
\end{proof}

\begin{theorem}[A high-content conditional class with no norm-support bound]
\label{thm:ep-content-class}
Fix $B\ge0$ and $\varepsilon>0$. There is an effective integer
$G_{B,\varepsilon}$ such that every primitive positive triple with
\[
 g\ge G_{B,\varepsilon},\qquad
 W_{g^{1/6}}(T)\le g^B
\]
satisfies $c\le R^{1+\varepsilon}$. In particular the stronger condition
$E_3^{>g^{1/6}}(T)\le g^B$ suffices. The large-prime hypothesis is not
asserted for all triples. In particular this is not a proof of standard $abc$.
\end{theorem}
\begin{proof}
The small-prime contribution to $\log(E_3/C_3)$ is at most the left side of
\eqref{eq:ep-content-small}; the large-prime contribution is at most
$B\log g$. Since $t\ge g\log7/2$, the sum of $\log8$, the angular penalty,
and $\log(E_3/C_3)$ divided by $t$ is at most
\[
 \frac{2\log8}{g\log7}
 +2C^2\frac{\log(3+g)}g
 +\frac{C^2}{3}\frac{\log(3+g)\log g}{\sqrt g}
 +\frac{2B\log g}{g\log7}.
\]
This tends to zero independently of the norm primes. Choose the threshold
so that it is at most $3\varepsilon/(1+\varepsilon)$ and apply
\eqref{eq:ep-height} after taking logarithms.
\end{proof}

There is also a usual uniform-constant formulation on any explicitly fixed
class for which $g\ge f(\log c)$ with a fixed function $f(t)\to\infty$ and
the same excess hypothesis: beyond a threshold Theorem~\ref{thm:ep-content-class}
applies, and below it finitely many triples are absorbed into $K_\varepsilon$.
No assertion that $g$ diverges on all prospective counterexamples is made.

\begin{theorem}[Several blocks, including global content one]
\label{thm:ep-few}
Fix an integer $d\ge1$ and reals $0<\delta<1$, $B\ge0$.
For every $\varepsilon>0$ there is a constant
$K_{d,\delta,B,\varepsilon}>0$ such that the standard $1+\varepsilon$
inequality holds for every primitive positive triple admitting a partition with
\begin{equation}\label{eq:ep-profile-condition}
 m\le d,\qquad \prod_{j=1}^m k_j\ge t^{m-1+\delta},\qquad
 W_{t^{\delta/12}}(T)\le t^B.
\end{equation}
These conditions define the class for $t\ge2$; the finitely many smaller
triples may be included by enlarging the same constant. No restrictions on
individual $q_i$ or their total count are imposed.
\end{theorem}
\begin{proof}
Weighted AM--GM in \eqref{eq:ep-regroup} gives
\[
 \prod_j\log Q_j\le\frac{(2t/m)^m}{\prod_j k_j},\qquad
 S\le\frac{2t}{\log7}.
\]
For $m\le d$, Theorem~\ref{thm:ep-main} therefore gives
\[
 \mathcal B_{\mathcal P}\le C_d t^{1-\delta}\log(3+t)
 \le t^{1-\delta/2}
\]
past a threshold depending only on $d,\delta$. With $Y=t^{\delta/12}$,
\eqref{eq:ep-small} is at most
$(\delta/12)t^{1-\delta/4}\log t=o(t)$ uniformly. The small-prime contribution to $\log(E_3/C_3)$ is no greater than
this full small-prime bound; the high-prime signed cost is at most
$B\log t=o(t)$. In \eqref{eq:ep-height} absorb the total cost
below $3\varepsilon t/(1+\varepsilon)$; this proves the result above a
uniform threshold. Below it $K=\max(1,c_0)$ suffices.
\end{proof}

\subsection{Actual profiles not covered by a polylogarithmic norm-prime cutoff}
There are infinitely many primes $q\equiv1\pmod3$. An elementary proof is to
suppose a finite list, take $x$ equal to three times their product, and take a
prime divisor of $x^2+x+1$. It is neither three nor any listed prime, and the
order of $x$ modulo it is three, so it is another prime congruent to one.
Choose $\pi_q$ of norm $q$ in the same ring, with one fixed orientation.

First take $k=\lceil(\log q)^2\rceil$ and $z=\pi_q^k$. Its coordinates are
primitive; a common rational prime divisor would require a conjugate factor.
Its boundary is nonzero because its norm is greater than one and a primitive
zero-boundary pair is a unit. After rotation it gives an actual triple with
\[
 g=k,\qquad \log c=\tfrac{k}{2}\log q+O(1)
             \asymp(\log q)^3.
\]
Thus $q>(\log c)^A$ eventually for every fixed $A$ while the one-block
angular and small-prime bounds apply uniformly. This demonstrates an infinite
extension of the \emph{profile} domain, not an infinite family satisfying the
additional unproved high-prime excess hypothesis.

For a content-one example, let $q\ne7$, put $k=\lceil(\log q)^3\rceil$, and
use $z=\pi_q^k(2+\zeta)^{k+1}$. It is again primitive and nondegenerate,
its global content is $\gcd(k,k+1)=1$, and the two indicated blocks satisfy
\[
 \log c\asymp(\log q)^4,\qquad k(k+1)\asymp(\log q)^6.
\]
Hence their product exceeds $t^{1+\delta}$ for any fixed $\delta<1/2$ and
large $q$. Again $q$ exceeds every fixed power of $t$. The profile hypothesis
in Theorem~\ref{thm:ep-few} is consequently not just the earlier bounded-prime
condition or a disguised common-power condition.

\begin{proposition}[Exact limit on squarefree split-norm profiles]
\label{prop:ep-squarefree}
Suppose all split exponents $e_i$ equal one. Then the minimum of
\eqref{eq:ep-penalty} over all partitions is exactly
\[
 \min_{\mathcal P}\mathcal B_{\mathcal P}(z)
   =C^2\log4\,(\log M-e_0\log3).
\]
It is attained by the single block. In particular regrouping removes the
previous artificial factor $\log\log c$ at a prime-norm endpoint, but
cannot make this particular bound $o(\log c)$ there.
\end{proposition}
\begin{proof}
Every block has $k_j=1$, hence $S=m$ and
$\sum_j\log Q_j=\log M-e_0\log3=:V$. Set $l=\log7>1$.
Since every $\log Q_j/l\ge1$, their product is at least their average;
therefore $\prod_j\log Q_j\ge l^{m-1}V/m$.
Also $C^{m-1}l^{m-1}\ge2^{m-1}\ge m$ for $m\ge1$.
Thus every partition has penalty at least $C^2\log4\,V$, with equality
for the single block. Since $V=2\log c+O(1)$, the last assertion follows.
This concerns the proved bound, not a lower bound on the true angular defect.
\end{proof}

\section{A large-prime excess bound that cannot hold globally}
\label{sec:ep-obstruction}
An attempted global proof must not simply assert that every large-prime
excess factor is small. That assertion is false even on triples that already
satisfy the stronger bound $c<R$.

We use only the classical Linnik theorem: there are absolute $A_0,L>0$ such
that each reduced progression $a\bmod m$ contains a prime at most
$A_0m^L$. Its quantitative form is recorded, for example, in
Xylouris \cite{EPLinnik}. We do not require the best numerical exponent.

\begin{theorem}[Deep large-prime excess with negative abc defect]
\label{thm:ep-safe}
Fix an integer $h\ge4$. For each odd prime $p$ let $\ell_p$ be the least
prime in the progression $p^h-1\pmod {p^{h+1}}$, and put
\[
 (a_p,b_p,c_p)=(1,\ell_p,\ell_p+1).
\]
Then $v_p(c_p)=h$ and
\[
 p^h\le c_p\le (A_0+1)p^{(h+1)L'}
\]
for an absolute $L'\ge\max(1,L)$, after harmless enlargement of $A_0$.
Moreover
\[
 R_p=\rad(a_pb_pc_p)\ge2\ell_p>c_p.
\]
For every fixed $A>0$, $p>(\log c_p)^A$ eventually and
\begin{equation}\label{eq:ep-excess-no-go}
 \liminf_{p\to\infty}
 \frac{\log E_3^{>(\log c_p)^A}(a_pb_pc_p)}{\log c_p}
 \ge\frac{h-3}{(h+1)L'}>0.
\end{equation}
Thus neither a uniform polylogarithmic bound on this raw excess nor its
uniform $o(\log c)$ logarithmic bound is true on all primitive triples.
The family is not a counterexample to $abc$ and does not refute the restricted
norm-profile theorems above or in the preceding manuscript.
\end{theorem}
\begin{proof}
The residue $p^h-1$ is coprime to $p^{h+1}$. Linnik supplies the upper bound.
Write $\ell_p=p^h-1+jp^{h+1}$ with $j\ge0$. Then
$c_p=p^h(1+jp)$, which proves exact valuation $h$ and the lower bound.
The prime $\ell_p$ is larger than two, so it is odd and $2\mid c_p$;
also $\gcd(\ell_p,c_p)=1$. These distinct supported primes imply
$R_p\ge2\ell_p>c_p$.

The height upper bound gives $\log c_p\le\log(A_0+1)+(h+1)L'\log p$.
Consequently $p>(\log c_p)^A$ for large $p$. Its contribution to the
restricted excess is $p^{h-3}$. Divide the logarithm of this contribution by
the height bound and take the lower limit, proving \eqref{eq:ep-excess-no-go}.
Finally $c_p\ge p^h$ makes the family height-unbounded.
\end{proof}

\begin{proposition}[Explicit compensation and an angular obstruction]
\label{prop:ep-compensation}
On the family in Theorem~\ref{thm:ep-safe},
\[
 E_3(T)=E_3(c_p),\qquad C_3(T)=\ell_p^2C_3(c_p),\qquad
 \frac{E_3(T)}{C_3(T)}\le\frac{c_p}{\ell_p^2}.
\]
For every triple $1+b=c$, including this family,
\[
 -\log|(z/\bar z)^3-1|=\log c+O(1)
 \quad(c\to\infty),
\]
with an absolute implied constant. Therefore a uniform sublinear angular
penalty on \emph{all} triples is also impossible.
\end{proposition}
\begin{proof}
The prime $\ell_p$ occurs with exponent one in $T=\ell_pc_p$ and is absent
from $c_p$, giving the first two equalities. Use $E_3(c_p)\le c_p$ and
$C_3(c_p)\ge1$ for the inequality.
For $a=1,b=c-1$, $M=c^2-c+1$ and \eqref{eq:ep-cubic} gives exactly
\[
 |(z/\bar z)^3-1|=
 \frac{3\sqrt3\,c(c-1)}{(c^2-c+1)^{3/2}}.
\]
After multiplication by $c$, this tends to $3\sqrt3$, proving the claim.
\end{proof}

This identifies an obstruction to bounding the archimedean and positive
multiplicity losses independently, not to correlated methods. Indeed the
exact identity with $\theta=4ab/c^2$ is
\[
 \left(\frac cR\right)^3
 =\frac4\theta\frac{E_3(T)}{C_3(T)}.
\]
Proving the corresponding global correlated bound without new arithmetic
input would be another expression of $abc$, not a completed bridge theorem.

\section{Dependency audit, independent computations, and precise remaining gates}
\label{sec:ep-audit}
The positive proof chain is
\[
\begin{gathered}
 \text{actual prime-norm factors}\ \longrightarrow\
 \text{integer exponent blocks}\\
 \mathrel{\phantom{\longrightarrow}}\Downarrow\quad
 \text{two established logarithmic-form bounds}\\
 \text{uniform block penalty and small-prime bound}.
\end{gathered}
\]
The uniform $1+\varepsilon$ consequences additionally require their explicitly
stated high-prime excess conditions. The Linnik construction independently
proves that an unrestricted positive-excess smallness claim is false.
The prime-norm terminal case has just one split exponent equal to one: there
is no nontrivial exponent block to extract. It remains unresolved; the present
method does not supply a hidden small predecessor for it.

The exact verifier factors bounded norms, reconstructs their oriented ring
factors, compares exponent-block products to the original coordinates, and
independently checks norm products and primitive coordinates. It also certifies
finite prime-neighbour examples by trial division and checks their exact
valuation, radical and credit. Finite computations neither prove Linnik nor
estimate unknown logarithmic-form constants. Machine-readable bounds and counts
are recorded only after execution.

The Lean companion proves the finite commutative multiplication and norm
regrouping cores, universal arithmetic progression divisibility and
nondivisibility, and explicit integer inequalities. These statements do not
formalize the external logarithmic-form estimates, the general splitting
classification, the Linnik existence theorem, or all asymptotic absorption.
A successful scoped compile must not be reported as a full-repository build,
independent peer review, autonomous-subagent audit, or a proof of standard $abc$.

The retained gates are: control the \emph{correlated} large-prime multiplicity
and low-depth credit; handle triples with bounded block contents or many
arithmetically incompatible contents; and supply a new input for prime-norm
terminal cases. The original IUT all-place comparison, FCRT arithmetic residual,
Pell/Mersenne first-depth estimates, and geometric uniformity routes remain
open at their previously recorded statements. The universal upper bounds on
raw high-prime excess and on the angular defect alone are now specifically
refuted, while their parent correlated routes remain active.
